% \newpage

\section{Related work}
\label{sec:related_work}

\paragraph{Why Does Generation Diversity Matter?}
Language models often collapse toward a narrow set of high-probability responses, resulting in an ``Artificial Hivemind'' characterized by substantial intra- and inter-model homogeneity on open-ended prompts \cite{jiang2025artificialhivemindopenendedhomogeneity}. Zhang et al.~\cite{Zhang2025} further show that post-training alignment reduces diversity through biases in preference data, limiting performance in domains that require diverse outputs.
In scientific reasoning, automated research agents benefit from a wider exploration breadth \citep{zou2026fmlbenchbenchmarkingmachinelearning}, and mode collapse causes reduced exploration in solution spaces \citep{yuan2026agenticredevolvingagenticsystems}. In creative tasks, reduced diversity in LLM outputs also diminishes the creativity of LLM-assisted writing \citep{abdulhai2026llmsdistortwrittenlanguage,Anderson2024HomogenizationEO,padmakumar2023does}. In mathematical proofs, strategy diversity is critical for exploring distinct reasoning paths \citep{Wu2025, Cao2025TowardsAM}. These findings motivate methods that preserve quality while expanding the range of model outputs.

% Language models often collapse toward a narrow set of high-probability responses. Jiang et al.~\cite{jiang2025artificialhivemindopenendedhomogeneity} characterize this as an ``Artificial Hivemind,'' demonstrating substantial intra- and inter-model homogeneity on open-ended prompts, while Zhang et al.~\cite{Zhang2025} show that post-training alignment further amplifies this effect through biases in preference data. This lack of diversity is particularly problematic in domains requiring broad exploration. Diverse outputs improve exploration in scientific reasoning and research agents \citep{Zou2025, yuan2026agenticredevolvingagenticsystems}, enhance creativity in LLM-assisted writing \citep{abdulhai2026llmsdistortwrittenlanguage,Anderson_2024,padmakumar2023does}, and increase the effectiveness of repeated sampling for mathematical reasoning by exploring distinct proof strategies \citep{Wu2025, Cao2025TowardsAM}. These findings motivate methods that preserve quality while expanding the diversity of model outputs.

\paragraph{Measuring Diversity of LLMs.}
Diversity of LLM-generated content has been measured along several dimensions. Lexical metrics include distinct n-grams \citep{ippolito2019comparison}, Self-BLEU \citep{zhu2018texygen}, Measure of Textual Lexical Diversity (MTLD) \citep{McCarthy2010MTLDVA}, and compression-based homogeneity scores \citep{shaib2024standardizing}. Semantic metrics build on sentence embeddings \citep{reimers2019sentencebertsentenceembeddingsusing, wieting2018paranmt} and information-theoretic constructions such as the Vendi Score \citep{friedman2023vendiscorediversityevaluation} and its conditional variant \citep{jalali2024conditionalvendiscoreinformationtheoretic}, as well as gradient-based measures like G-Vendi \citep{jung2025prismaticsynthesisgradientbaseddata}. A complementary line of work studies how diversity collapses through the training pipeline \citep{kirk2024understandingeffectsrlhfllm, omahony2024attributingmodecollapsefinetuning, padmakumar2023does, west2025base, dang2025assessing}, and recent benchmarks evaluate humanlike or effective semantic diversity directly \citep{zhang2025noveltybenchevaluatinglanguagemodels, shypula2025evaluatingdiversityqualityllm, guo2025benchmarking, Shahid2025}. However, most evaluations assess diversity or quality in isolation, whereas we jointly evaluate capability and diversity across real-world applications to ensure models preserve multiple plausible answers without mode collapse \citep{Misaki2026,Puri2026} or catastrophic forgetting \citep{luo2025empiricalstudycatastrophicforgetting}.



% \liwei{Can we tighten the following paragraph a bit more? It's a bit too long.}
\paragraph{Improving Diversity during Training and Inference.}
Recent work has explored improving diversity through prompting \citep{Zhang2025, lau2025dipperdiversitypromptsproducing, lagzian2025multinoveltyimprovediversitynovelty, wang2025multilingualpromptingimprovingllm, Kim2026, Wu2025}, decoding \citep{holtzman2020curiouscaseneuraltext, nguyen2024turningheatminpsampling, ruan-etal-2025-g2, wang2024lpoadaptivedecodinglatent}, and training \citep{li2025jointlyreinforcingdiversityquality, lanchantin2025diverse, chung2025modifyinglargelanguagemodel, Chen2025PosttrainingLL, slocum2025diversepreferencelearningcapabilities, Puri2026} (see Appendix~\ref{appx:related-work-table} for a comparison table). At inference time, Verbalized Sampling \citep{Zhang2025} prompts models to express a distribution over possible responses, while String Seed of Thought \citep{Misaki2026} injects random strings as entropy sources. DIPPER \citep{lau2025dipperdiversitypromptsproducing} and Multi-Novelty \citep{lagzian2025multinoveltyimprovediversitynovelty} construct diverse prompt ensembles, and multilingual, persona, or chain-of-thought prompting can elicit variations \citep{wang2025multilingualpromptingimprovingllm, Wu2025}. At decoding time, min-$p$ sampling \citep{nguyen2024turningheatminpsampling} adaptively truncates low-probability tokens based on model confidence, while G2 \citep{ruan-etal-2025-g2} and adaptive temperature methods \citep{wang2024lpoadaptivedecodinglatent} guide generation with auxiliary diversity modules or learned per-instance decoding parameters. In post-training, RL has emerged as a powerful method for enhancing diversity. Multi-answer RL \citep{Puri2026} trains models to output multiple diverse answers within one response; Soft Preference Learning decouples entropy from KL-penalty to improve lexical and semantic variety \citep{slocum2025diversepreferencelearningcapabilities}; DivPO \citep{lanchantin2025diverse} applies DPO to responses that are both diverse and high quality. 

Several recent works share our goal of jointly optimizing quality and diversity through RL. DARLING \citep{li2025jointlyreinforcingdiversityquality} balances quality and diversity by rewarding a multiplicative aggregation of quality and diversity metrics; DQO \citep{Chen2025PosttrainingLL} uses a determinant-based group diversity reward, but assigns the same reward to all responses in the group. GRPO-Unlikeliness \citep{he2025rewardingunlikelyliftinggrpo} and GAPO \citep{anschel2025group} both encourage diversity through frequency- and likelihood-based rewards. However, although these methods explore different reward designs, none of them allow a single model to switch between producing one high-quality answer and a diverse set of responses.
\method, on the other hand, approaches quality-diversity alignment from a MARL-inspired perspective, treating diverse generation as a competition-and-coordination problem among role-conditioned agents within a shared policy. This formulation enables explicit per-role credit assignment under quality constraints, encouraging specialization across the high-quality response space while keep the model's standard behavior intact. 